% open-logic.tex
% Compilation de l'ensemble du contenu avec le style open-logic-dev

\documentclass[../include/open-logic-part]{subfiles}

\begin{document}

\clearpage

\begin{editorial}
Ce fichier est le pilote de compilation générale du projet Open
Logic. Il charge les parties de la collection sans sélection pour
un cours particulier.

Par défaut, tous les opérateurs logiques sont traités comme primitifs
et les démonstrations traitent tous les cas correspondants. Pour un
cours ou pour l'étude individuelle, on peut choisir les chapitres et
laisser certains cas en exercice.

Pour consulter des textes préparés pour l'enseignement ou l'étude,
voir les \href{http://builds.openlogicproject.org/}{exemples de cours
  proposés sur le site du projet}. Pour composer votre propre texte,
vous pouvez partir du
\href{https://github.com/OpenLogicProject/OpenLogic/tree/master/courses/sample}{fichier
  pilote d'exemple} ou examiner les sources des manuels dérivés, qui
offrent des exemples plus élaborés.
\end{editorial}

\olimport[sets-functions-relations]{sets-functions-relations-complete}

\olimport[propositional-logic]{propositional-logic}

\olimport[first-order-logic]{first-order-logic}

\olimport[model-theory]{model-theory}

\olimport[computability]{computability}

\olimport[turing-machines]{turing-machines}

\olimport[incompleteness]{incompleteness}

\olimport[second-order-logic]{second-order-logic}

\olimport[lambda-calculus]{lambda-calculus}

\olimport[many-valued-logic]{many-valued-logic}

\olimport[normal-modal-logic]{normal-modal-logic}

\olimport[applied-modal-logic]{applied-modal-logic}

\olimport[intuitionistic-logic]{intuitionistic-logic}

\olimport[counterfactuals]{counterfactuals}

\olimport[set-theory]{set-theory}

\olimport[methods]{methods}

\olimport[history]{history}

\olimport[reference]{reference}

\end{document}
